\section{The original formal arcs and their specialization}
\label{sec:pel-formal-arcs}

The remaining formal-arc clause prescribes the period factors, the
torsion labels and the specialization map exactly.  Its input
already says that every original period factor is a unit.  Thus
the construction has one monodromy ray, including when the
coefficient ring is itself a discrete valuation ring.

\begin{theorem}[The exact original arc clause]
\label{thm:pel-original-arcs}
Every canonical local model constructed above satisfies the full
source clause \texttt{N4}: the prescribed power-series arc exists
uniquely, has exactly the prescribed abelian locus, realizes the
original signed period and torsion formulas, and has the required
specialization for every ring map in that clause.
\end{theorem}

Let $O'$ and $a$ be the source's \texttt{ArcIn} data.  Thus $O'$
is a characteristic-zero field, with $\varpi=1$, or a complete
discrete valuation ring with uniformizer $\varpi$.
Write $B_0,\mathscr M_0,X,c_0,\psi^0,\tau^0$ for its original
abelian and period data.  The nonzero lattice $X$ has the original
$\mathcal O_E$-action.  The specified symmetric integral form
$b$ is positive definite, adjoint-compatible, and takes values
in $2M'\mathbb Z$, where $M'=\operatorname{lcm}(M,A)$.
Put
\[
       R=O'[[t]],\quad I=(t),\quad Q(y)=b(y,y)/2.
\]
No division in $O'$ is intended: $Q(y)$ is an integer by the
given divisibility, including in residue characteristic two.

\subsection*{The actual period construction}

Keep $B=B_0\times_{O'}R$, its original line, rigidification,
coefficient action and point map $c$.  Change only the prescribed
period generators, by
\[
       \psi_y=\psi^0_y\,t^{Q(y)},\qquad
       \tau_{y,z}=\tau^0_{y,z}\,t^{b(y,z)}.
\]
The second formula follows from the first and the original
definition of $\tau$, since
$Q(y+z)-Q(y)-Q(z)=b(y,z)$.
Every $\psi^0_y$ is an integral unit generator by
\texttt{ArcIn.gen}.  Positivity of $Q$ gives
$\psi_y\in I$ for $y\ne0$, and for every $n$ all but finitely
many $\psi_y$ lie in $I^n$.  The original evenness and
bimultiplicativity are preserved.  These are the complete
\texttt{DegData} conditions for the ideal $(t)$, not merely for
the maximal ideal of $R$.

The coefficient adjunction identity for $b$ preserves the
coefficient action on the new period data.  More specifically,
with $w=w_0$ and $w+\bar w=1$, symmetry and adjunction give
\[
               b(y,wy)=b(y,\bar w y)=Q(y).
\]
The original normalized generator $u$ in
\texttt{TraceFormTwistIso} consequently still satisfies
$\tau_{y,wy}u(c(y))=\psi_y\,\mathrm{rig}$.
Its line, with exponents $(1,-2,-1,2)$, is precisely the
primitive-line comparison on $B$.  The last identity says that
this comparison respects the new period action.

Apply the actual Mumford degeneration equivalence over the
complete noetherian normal pair $(R,(t))$.  It constructs a
semiabelian scheme $G/R$ and its actual symmetric ample cubical
line $\mathscr L$.  The period lattice is the same $X$, with
$\phi=1$.  The degree formula gives a principal generic
polarization.  Tensor compatibility carries the original
coefficient, adjunction and primitive identities to this family,
so it is the required original \texttt{SAB}, with no extra
power or scalar\cite[Theorems 4.2.1.14 and 4.4.16]{formal-arcs-lan}.

The integral expansion and all-integer torsion arguments of
Theorem~\ref{thm:dvr-original-pel-unif} work for this complete
pair: their formal comparison is on every $R/(t^{e+1})$, and
the needed period positivity has just been proved directly.
They supply the exact \texttt{PELUnif}, the original line
comparison and the theta-action description of all torsion.

\subsection*{The original finite labels}

Let $k=M$ or $k=A$, and let $(y,z)$ be one of the original
selected \texttt{QkPt}s over $\operatorname{Frac}O'$.
Keep its point on $B_0$ and change its inverse-biextension
coordinates by
\[
                       z_t.s(\chi)
          =z.s(\chi)\,t^{b(y,\chi)/k}.
\]
The exponent is an integer because $k\mid M'$.
Bilinearity makes this a point of the same torus extension, and
the full equality $[k]z=\iota^0(y)$ becomes
$[k]z_t=\iota(y)$.  Period equivalence is preserved because
\[
     b(y+kx,\chi)/k-b(y,\chi)/k=b(x,\chi).
\]
Addition and the coefficient action are preserved too, using
$b(ay,\chi)=b(y,\bar a\chi)$.
The change in the torsion commutator is
$t^{(b(y,y')-b(y',y))/k}=1$; thus the original Weil-pairing
root and sign are unchanged.

The source's geometric full-Tate condition also persists.
Over any algebraically closed generic field choose compatible
roots $t^{1/n}$ for all $n$, and apply the same formula to the
original \texttt{TateIsoQ}, using
$(t^{1/n})^{b(y,\chi)}$.  The period, addition and tower
identities above give the entire coherent Tate marking, and
symmetry preserves its pairing roots.  For $n=M,A$, divisibility
reduces it to the displayed integral-exponent formula, independent
of the chosen roots.  The original Lie multiplicities on $B_0$,
together with the constant torus contribution
$\operatorname{rank}_{\mathcal O_E}X$ at each embedding, give
exactly the original signature.  This retains the entire period
rank and the source's full-Tate, rather than merely finite-label,
hypothesis.

Only the prescribed restriction of $a.\eta$ is used here.
Choose $e_+\in\overline{\mathfrak b}$ with
$1-e_+\in\mathfrak b$, and define the new half-label at
$x\in\Lambda$ from the transformed point $a.\eta(e_+x)$.
The geometric Tate comparison on
$\overline{\mathfrak b}\Lambda$ proves that this is independent
of the choice, is $\mathcal O_E$-linear, has kernel
$\mathfrak b\Lambda$, and agrees with the prescribed transformed
points on $\overline{\mathfrak b}\Lambda$.
Indeed $\mathfrak b\overline{\mathfrak b}=M\mathcal O_E$,
so the changes in these assertions lie in $M\Lambda$.
Geometric equality descends for these rational torsion points.
The full-$A$ label uses the given $a.\eta_A$ on all of $\Lambda$.
Thus no additivity of the unconstrained values of $a.\eta$
outside the selected submodule is assumed.

Consequently $G$ gives the original moduli object over
$D(\varpi t)$.  The label coordinates are defined there over
the original coefficient field: their only new factors are
the displayed integral powers of $t$.  The actual torsion
quotient is defined over that field, so no geometric point or
coefficient extension has been substituted for a rational label.

\subsection*{The map to the prescribed toroidal model}

The original period factors are unit generators on all of $O'$.
Hence, for every discrete valuation centered on
$\operatorname{Spec}R$, the monodromy pairing is
\[
                           v(t)\,b.
\]
There is no additional $\varpi$-direction contribution.
The cusp lattice, coefficient action and admissible filtration
are obtained from the same exact all-level torsion sequence and
the transported full Tate marking.  In particular the torus
submodule and its annihilator, and hence the filtration and
torus argument at the original level, are fixed as the centered
valuation varies at a point of $t=0$.  On the original singleton
lattice the auxiliary Siegel degeneration is precisely this
family with the verified finite-level symplectic relabeling;
its period form is the corresponding image of $b$.
For any additional auxiliary presentation the induced tuple
has the images of this same filtration and form, as in the
proof of Lan's Theorem 11.4.  Thus the pairings lie in one ray
in the \emph{specified} cusp cone, including every member of
the auxiliary tuple.  The fixed admissible cone decomposition
contains that ray in a cone; at points where $t$ is a unit the
pairing is zero.  This is exactly the centered-valuation cone
condition of Theorem 7.14, with the original filtration and
torus argument retained.  The exact toroidal universal property gives
\[
              \chi:\operatorname{Spec}O'[[t]]\longrightarrow Y,
\]
extending the original moduli object and carrying its actual
family\cite[Theorem 11.4]{formal-arcs-ramified}.
This applies to the possibly two-dimensional normal base itself;
separate extensions along curves would not suffice.

The abelian locus is exactly $D(t)$.  Indeed, the subgroup of
unit periods at a point is all of $X$ if $t$ is a unit and is
zero otherwise, since $b(y,y)>0$ for $y\ne0$.  The exact
unit-period torsion sequence then gives torus rank zero in the
first case and $\operatorname{rank}_{\mathbb Z}X>0$ in the
second\cite[Proposition 4.5.3.11 and Corollary 4.5.3.12]
{formal-arcs-lan}.  The map on $D(\varpi t)$ lies in the
original generic moduli interior, as already constructed.

With the identity base change of $B_0$, its original line and
the identity on $X$, the period and torsion formulas above are
exactly the two \texttt{GenRel} identities in
\texttt{HasArcData}: the factors are
$t^{b(y,y)/2}$ and $t^{b(y,z)/k}$, with their positive signs.
The previously verified \texttt{PhiRel} calculation identifies
the latter points with the actual original labels, for every
geometric field and every lift appearing in that definition.
This proves the asserted existence with all its stated fields.

\subsection*{Why the source's expansion contract gives uniqueness}

We record the needed recovery argument, since a vector-space
bijection of theta functions alone would not determine the
abelian variety.  Suppose two source \texttt{PELUnif}s have
the same integral abelian, line and period data, up to the
specified rigidified isomorphisms.  Their all-power
\texttt{ExpBij}s identify the graded section spaces
\[
                     \bigoplus_{k\geq0}H^0(A,\mathscr L^k).
\]
In degree zero use the fraction field as the scalar ring.
This identification is an isomorphism of graded rings.
To check multiplication, clear the two actual common
denominators in \texttt{Expands}.  On every integral thickening,
\texttt{ExpandsAt} is the literal finite character expansion.
The source's \texttt{t\_mul} and the cube identity multiply it
by coefficient convolution.  The coefficient-line products
are those of the same original $B$ and its biextension.
The compatibility fields \texttt{gt\_t} and
\texttt{w\_compat} make these products agree on all thickenings.
Convolution is finite on each thickening by \texttt{ExpandsAt};
the compatible coefficient limits are actual sections, because
the coefficient lines lie on the proper abelian part and their
finite section modules are complete over $R$.
Thus after clearing the two common denominators we have an
actual coefficient family, before using \texttt{ExpBij}.
Separatedness of the complete ring and the nonzero common
denominators give the equality of generic products.  Thus no
ring structure is inferred solely from equal dimensions.

The generic line is ample.  The entire graded section ring
therefore recovers the polarized generic variety by
$\operatorname{Proj}$, with its original line.
The resulting isomorphism preserves the identity: evaluation
at zero is the same convergent coefficient sum, because the
source's \texttt{t\_unit} is one and \texttt{w\_rig}
preserves the original rigidification.  Clear a denominator
and check this evaluation on every thickening as before.
An isomorphism of abelian varieties preserving zero is a group
isomorphism.

It also preserves the original coefficient action.  Extend
the generic group isomorphism and its inverse over the normal
base by semiabelian full faithfulness, and extend the
rigidified cubical line comparison.  Its formal restriction
is the one prescribed by the common torus and abelian
coordinates: pullbacks of every integral theta power agree
on all thickenings.  Here the line on \emph{each} competing
source \texttt{SAB} is relatively ample: it is symmetric
cubical, and its generic fibre is ample by that source
\texttt{PELUnif}'s polarized abelian variety; Lan's
Lemma 4.2.1.6 applies over the normal base.
The relatively ample line therefore
separates maps through its canonical relative-Proj open
immersion\cite[Lemma 4.2.1.6 and Propositions 3.3.1.5,
3.3.2.2]{formal-arcs-lan}\cite{formal-arcs-relative-proj}.
The \texttt{OLinear} fields identify the
coefficient actions in those formal coordinates: they give the
same map on the abelian quotient and on the torus.  Two such
endomorphisms of a torus extension differ by a homomorphism
from its proper abelian quotient to the torus; this homomorphism
is zero.  Their
composites consequently agree formally, and then generically
by separatedness and the integral thickening comparisons.
The extended group isomorphism intertwines them.

Finally the \texttt{PhiRel} condition determines each selected
torsion label.  Under the recovered line and group
isomorphism it specifies the same projective theta action.
Two $k$-torsion points with that action have quotient pairing
trivially with every $k$-torsion point.  Perfectness of the
principal Weil pairing in characteristic zero makes them
equal.  For the half-label it suffices to check the source's
selected arguments in $\overline{\mathfrak b}\Lambda$:
$\Lambda=\mathfrak b\Lambda+
\overline{\mathfrak b}\Lambda$, and the original label
vanishes on $\mathfrak b\Lambda$.  Thus the entire original
half-label and full-$A$ label maps are recovered.

Now take any competing $\chi'$ in the source's proposed
characterization.  Its \texttt{HasArcData} identifies its
abelian part, rigidified line, coefficient action, period
points and factors with precisely the twisted data above.
The normalized line comparison preserves the cube and
symmetry, which are unique on the abelian part.
The recovery argument identifies the generic original
moduli objects and their finite labels.  The fine moduli
property makes the generic points equal, and separatedness
of $Y$ gives $\chi'=\chi$ on all of $\operatorname{Spec}R$.
Conversely the constructed $\chi$ has every stated property.
This proves the exact ``if and only if'' uniqueness clause.

\subsection*{The prescribed specialization}

Suppose $O'$ is a complete discrete valuation ring.  Evaluation
at $t=\varpi$ is defined by the convergent power series and is
compatible with the adic pairs $(R,(t))\to(O',(\varpi))$.
In the Mumford construction this base change gives the
degeneration data
\[
                  \psi^0_y\,\varpi^{Q(y)},
                  \qquad \tau^0_{y,z}\,\varpi^{b(y,z)},
\]
and the same label formulas with $t$ replaced by $\varpi$.
Because this base change is not flat, we verify it in the
actual construction.  Choose the same sufficiently positive
auxiliary powers and finite character set for the relatively
complete Mumford model.  The required positivity inequalities
depend only on the integer exponents of $b$, since all the
original period coefficients are units, so the same choice
works before and after evaluation.

Lan's affine charts are character-graded algebras generated by
the coefficient modules
\[
 I_y^{-1}I_{y,\alpha}^{-1}I_zI_{z,\beta}\,
       \mathcal O_{\beta-\alpha+\phi(z-y)}.
\]
Here every $I$ is a unit times a principal integral power of
$t$, and the character line $\mathcal O_\gamma$ is pulled
back from the original abelian scheme over $O'$.
For a fixed character, the coefficient ideal is the sum of
the resulting powers of $t$.  If their exponents are bounded
below, this is the principal power with the smallest exponent;
if they are unbounded below, it is the localization $R[1/t]$.
Both descriptions commute with the specified evaluation:
$t^aR$ becomes $\varpi^aO'$, and $R[1/t]$ becomes $O'[1/\varpi]$.
Consequently each character module and its multiplication
base change to the identical module with $t^a$ replaced by
$\varpi^a$.  This remains true for the fractional modules
appearing on localized charts.  The torus coordinates and
the period translation maps specialize by the same explicit
formulas\cite[Corollary 4.5.1.8 and the following affine-chart
formula]{formal-arcs-lan}.

The formal quotient is constructed on each thickening by
gluing the translated charts and taking the represented
fpqc period quotient.  These operations and the descended
line commute with this chartwise base change.  Thus the
proper formal Mumford model, with its line, specializes to
the model for the displayed $\varpi$-periods.
Uniqueness of proper algebraization identifies their actual
proper models and semiabelian opens; the tensor-compatible
line construction recovers the original primitive line
from the auxiliary powers\cite[Proposition 4.5.2.15 and
Theorem 4.4.16]{formal-arcs-lan}.
In particular we have verified the period action, not just
the formal Raynaud group.  No fraction-field map sending the whole
$\operatorname{Frac}(O'[[t]])$ to $\operatorname{Frac}O'$
is needed or asserted.

If a map $y':\operatorname{Spec}O'\to Y$ has the source's
specified \texttt{HasArcData} and generic interior property,
the recovery argument over $(O',(\varpi))$ identifies its
generic moduli object with this specialization, including
all original labels.  Separatedness gives equality of the
two maps on $\operatorname{Spec}O'$.

The conclusion holds for every ring map $\mathrm{ev}$ in
\texttt{N4}, not only for a chosen continuous one.  If such a
map fixes $O'$ and sends $t$ to $\varpi$, the difference
between a series and its first $r$ terms maps into
$\varpi^rO'$.  Completeness and separatedness force its value
to be the ordinary evaluation sum.  Thus every permitted
$\mathrm{ev}$ is that evaluation map, and the exact final
specialization equality follows\cite{formal-arcs-source}.

Together with the preceding \texttt{N0}, \texttt{N2},
\texttt{SplitFibres}, \texttt{N3} and \texttt{N6} suppliers,
this establishes the original geometric \texttt{PartI}
construction.  Quantitative height, section-remainder,
boundary comparison and global theta-norm estimates remain
separate arithmetic obligations; no whole proof-map node is
declared complete by this construction alone.
